\section{Représentation des programmes}
\subsection{Limitations des AST}

La structure normalement utilisée par les algorithmes opérant sur des programmes
est l'\textit{arbre abstrait de syntaxe} ou AST (\textit{abstract syntax tree}
en anglais). Ceux-ci sont une représentation arborescente du code source
encodant la structure du programme. Malheureusement ils ont deux désavantages
majeurs pour notre propos :

\begin{itemize}
    \item les AST n'encodent pas l'ordre d'exécution du programme, celui-ci est
          défini à part en fonction des étiquettes des nœuds de l'arbre. Cela
          complique l'identification des opérations survenant après un appel
          récursif~;
    \item il est difficile de définir comment scinder un AST en deux, ce qui est
          nécessaire dans notre algorithme pour séparer les opérations venant
          après un appel récursif afin de les placer dans une continuation.
\end{itemize}

Pour ces raisons, nous utilisons une autre représentation des programmes qui
représente bien comment le programme est évalué.

\subsection{Formes linéaires}

\begin{definition}
    Une sous-expression $e'$ d'une expression \texttt{OCaml} $e$ est dite
    \textit{localement terminale} si l'évaluation de $e$ se termine
    immédiatement après l'évaluation de $e'$, à moins que $e'$ ne soit pas
    évaluée pour évaluer $e$.
\end{definition}

\begin{example}
    Dans l'expression \texttt{if }$a$\texttt{ then }$b$\texttt{ else }$c$, $a$
    n'est pas localement terminale car son évaluation est forcément suivie de
    l'évaluation de $b$ ou $c$. En revanche, $b$ et $c$ sont toutes les deux
    localement terminales car l'évaluation du \texttt{if} se termine
    immédiatement après l'évaluation de $b$ ou $c$.

    Dans l'expression $a$\texttt{ + }$b$, ni $a$ ni $b$ ne sont localement
    terminales car leur évaluation est toujours suivie de l'opération d'addition
    avant la fin de l'évaluation de $a$\texttt{ + }$b$.
\end{example}

\begin{definition} On définit par induction mutuelle les
    \textit{opérations élémentaires} et les \textit{formes linéaires}.


    Les \textit{opérations élémentaires} sont des entités syntaxiques définies
    par les mêmes règles que les expressions \texttt{OCaml}, mais légèrement
    modifiées :

    si $C(e_1, \ldots, e_n)$ est une règle de construction d'une expression
    \texttt{OCaml} avec $e_1, \ldots, e_n$ ses sous-expressions, alors la règle
    de construction correspondante pour les opérations élémentaires est \\
    $C(a_1, \ldots, a_k, l_1, \ldots, l_{n - k})$ avec $a_i$ des noms de
    variables \texttt{OCaml}, une pour chaque sous-expression
    \textbf{non localement terminale}, et $l_i$ des \textit{formes linéaires},
    une pour chaque sous-expression \textbf{localement terminale}.

    Une \textit{forme linéaire} est une liste de couples $(a, e)$ avec $a$ un
    nom de variable \texttt{OCaml} et $e$ une opération élémentaire.
\end{definition}

On notera la forme linéaire $[(a_1, e_1), (a_2, e_2), \ldots]$ comme suit :
\begin{align*}
    a_1 & := e_1 \\
    a_2 & := e_2 \\
    \ldots
\end{align*}

\begin{remark}
    Intuitivement, les formes linéaires listent les étapes élémentaires de
    l'évaluation, et le résultat de chaque opération est stocké dans le nom de
    variable associé. Le résultat de l'évaluation de la forme linéaire entière
    est stocké dans le nom de variable du dernier couple.
\end{remark}

\begin{example}
    On représente ici une expression \texttt{OCaml} et une forme linéaire
    équivalente:

    \small
    \begin{minipage}{0.3\textwidth}
        \begin{minted}{ocaml}
match l with
| [] -> 0
| x :: q -> x + sum q
        \end{minted}
    \end{minipage}%
    \begin{minipage}{0.7\textwidth}
        \begin{align*}
            a_1 := \  & \texttt{l}                                             \\
            a_2 := \  & \texttt{match } a_1 \texttt{ with}                     \\
                      & \texttt{| [] -> } a_3 := 0                             \\
                      & \texttt{| x :: q -> } \begin{aligned}[t]
                                                  a_4 := \  & \texttt{q}           \\
                                                  a_6 := \  & \texttt{sum}         \\
                                                  a_7 := \  & a_6 \ a_4            \\
                                                  a_8 := \  & \texttt{x}           \\
                                                  a_9 := \  & a_8 \texttt{ + } a_7
                                              \end{aligned}
        \end{align*}
    \end{minipage}

    \normalsize
    \bigskip

    Chaque ligne correspond à l'exécution d'une opération élémentaire : charger
    la valeur d'une variable, effectuer un appel de fonction, faire une addition...
    On a distingué ici les variables présentes dans le code original en
    \texttt{monospace} et les variables introduites pour créer la forme linéaire
    en \textit{italique} mais celles-là sont en réalité du même type.
\end{example}

\subsection{Linéarisation}

On peut définir un algorithme transformant une expression (représentée par un
AST) en une forme linéaire équivalente.

\begin{algorithm}[H]
    \caption{Linéariser}
    \Entree{$e$ une expression \texttt{OCaml}}
    \Sortie{$l$ une forme linéaire correspondant à $e$}

    \uSi{$e$ est une addition $a$\texttt{ + }$b$}{
        $l_a \leftarrow$ Linéariser($a$) \\
        $l_b \leftarrow$ Linéariser($b$) \\
        $a_i \leftarrow$ Le nom de variable du dernier couple de $l_a$ \\
        $a_j \leftarrow$ Le nom de variable du dernier couple de $l_b$ \\
        $a_k \leftarrow$ Un nom de variable non utilisé dans $e$, $l_a$ ou $l_b$ \\
        Renvoyer $l_a \cdot l_b \cdot (a_k, a_i\texttt{ + }a_j)$
    }\uSinon{
        \strut \ldots \ // On continue de même pour chaque règle de construction
        des expressions \texttt{OCaml}.
    }
\end{algorithm}

On peut définir similairement un algorithme Délinéariser transformant une forme
linéaire en un AST \texttt{OCaml} en remarquant que chaque instruction $a := e$
peut être représentée par la construction \texttt{let a = e in}. On peut aussi
démontrer par induction sur l'ensemble des expressions \texttt{OCaml} que les
algorithmes Linéariser et Délinéariser préservent le comportement du programme,
c'est-à-dire que les expressions $e$ et Délinéariser(Linéariser($e$)) sont
équivalentes.

\begin{remark}
    En pratique, nous utilisons un algorithme Délinéariser légèrement modifié
    pour améliorer la lisibilité des programmes en sortie, mais celui-ci ne
    change pas la correction de l'algorithme.
\end{remark}
